
\documentclass{article}
\usepackage{colortbl}
\usepackage{makecell}
\usepackage{multirow}
\usepackage{supertabular}

\begin{document}

\newcounter{utterance}

\centering \large Interaction Transcript for game `codenames', experiment `risk\_high', episode 2 with claude{-}opus{-}4{-}8{-}azure.
\vspace{24pt}

{ \footnotesize  \setcounter{utterance}{1}
\setlength{\tabcolsep}{0pt}
\begin{supertabular}{c@{$\;$}|p{.15\linewidth}@{}p{.15\linewidth}p{.15\linewidth}p{.15\linewidth}p{.15\linewidth}p{.15\linewidth}}
    \# & $\;$A & \multicolumn{4}{c}{Game Master} & $\;\:$B\\
    \hline

    \theutterance \stepcounter{utterance}  
    & & & \multicolumn{2}{p{0.3\linewidth}}{
        \cellcolor[rgb]{0.9,0.9,0.9}{
            \makecell[{{p{\linewidth}}}]{
                \texttt{\tiny{[GM$|$GM]}}
                {'hidden': {'team': ['ornitorinco', 'lettiera', 'germania', 'veleno', 'spazio', 'radice', 'canada', 'rete', 'tazza'], 'innocent': ['blocco', 'morte', 'aquila'], 'opponent': ['oliva', 'fantasma', 'strega', 'stadio', 'pupilla', 'torta', 'acero', 'carta'], 'assassin': ['bufalo', 'colpo', 'rivoluzione', 'berlino', 'talpa']}, 'revealed': {'team': {'team': [], 'innocent': [], 'opponent': [], 'assassin': []}, 'opponent': {'team': [], 'innocent': [], 'opponent': [], 'assassin': []}}}
            }
        }
    }
    & & \\ \\

    \theutterance \stepcounter{utterance}  
    & & \multicolumn{4}{p{0.6\linewidth}}{
        \cellcolor[rgb]{0.9,0.9,0.9}{
            \makecell[{{p{\linewidth}}}]{
                \texttt{\tiny{[P1$\langle$GM]}}
                \texttt{Fornisci un indizio di una sola parola che si riferisca ad alcune parole della tua squadra dall'elenco sottostante. Puoi scegliere a quante e quali parole si riferisce il tuo indizio.} \\
\texttt{L'obiettivo è far sì che il tuo partner indovini tutte le parole della tua squadra in base agli indizi che gli fornisci, prima che l'altra squadra abbia scoperto tutte le sue. Se il tuo partner indovina una parola 'avversaria' o 'distrattore', il suo turno di indovinare termina e la squadra avversaria può fare un turno. Se il tuo partner indovina una parola 'kill', il gioco termina e la tua squadra perde immediatamente.} \\
\texttt{Più sono le parole a cui si riferisce il tuo indizio, più velocemente puoi vincere il gioco, ma le associazioni possono essere più difficili da indovinare per il tuo partner, quindi scegli saggiamente la tua parola indizio. La parola indizio deve essere semanticamente correlata alle parole target, non può essere una delle parole delle liste o contenere parti di esse.} \\
\\ 
\texttt{Fornisci sempre il tuo indizio di una sola parola e l'elenco separato da virgole delle parole target correlate nel seguente formato e rendi le tue risposte il più brevi possibile, non includere mai altro testo oltre a quello richiesto in questo formato:} \\
\\ 
\texttt{INDIZIO: <PAROLA>} \\
\texttt{OBIETTIVI: <OBIETTIVI>} \\
\\ 
\texttt{Le parole della tua squadra sono: ornitorinco, lettiera, germania, veleno, spazio, radice, canada, rete, tazza.} \\
\texttt{Le parole della squadra avversaria sono: oliva, fantasma, strega, stadio, pupilla, torta, acero, carta.} \\
\texttt{Le parole distrattore sono: blocco, morte, aquila.} \\
\texttt{Le parole kill sono: bufalo, colpo, rivoluzione, berlino, talpa.} \\
            }
        }
    }
    & \\ \\

    \theutterance \stepcounter{utterance}  
    & \multicolumn{4}{p{0.6\linewidth}}{
        \cellcolor[rgb]{0.9,0.9,0.9}{
            \makecell[{{p{\linewidth}}}]{
                \texttt{\tiny{[P1$\rangle$GM]}}
                \texttt{INDIZIO: Nazione} \\
\texttt{OBIETTIVI: germania, canada} \\
            }
        }
    }
    & & \\ \\

    \theutterance \stepcounter{utterance}  
    & & & \multicolumn{2}{p{0.3\linewidth}}{
        \cellcolor[rgb]{0.9,0.9,0.9}{
            \makecell[{{p{\linewidth}}}]{
                \texttt{\tiny{[GM$|$GM]}}
                \texttt{nazione} \\
            }
        }
    }
    & & \\ \\

    \theutterance \stepcounter{utterance}  
    & & & \multicolumn{2}{p{0.3\linewidth}}{
        \cellcolor[rgb]{0.9,0.9,0.9}{
            \makecell[{{p{\linewidth}}}]{
                \texttt{\tiny{[GM$|$GM]}}
                ['germania', 'canada']
            }
        }
    }
    & & \\ \\

    \theutterance \stepcounter{utterance}  
    & & & \multicolumn{2}{p{0.3\linewidth}}{
        \cellcolor[rgb]{0.9,0.9,0.9}{
            \makecell[{{p{\linewidth}}}]{
                \texttt{\tiny{[GM$|$GM]}}
                {'word': 'germania', 'assignment': 'team'}
            }
        }
    }
    & & \\ \\

    \theutterance \stepcounter{utterance}  
    & & & \multicolumn{2}{p{0.3\linewidth}}{
        \cellcolor[rgb]{0.9,0.9,0.9}{
            \makecell[{{p{\linewidth}}}]{
                \texttt{\tiny{[GM$|$GM]}}
                {'word': 'canada', 'assignment': 'team'}
            }
        }
    }
    & & \\ \\

    \theutterance \stepcounter{utterance}  
    & & \multicolumn{4}{p{0.6\linewidth}}{
        \cellcolor[rgb]{0.9,0.9,0.9}{
            \makecell[{{p{\linewidth}}}]{
                \texttt{\tiny{[GM$\rangle$P2]}}
                \texttt{Fornisci un elenco separato da virgole di un massimo di 2 parole dal seguente elenco che meglio si riferiscono o sono più strettamente associate alla parola 'nazione'. Inizia sempre il tuo elenco di ipotesi con INDOVINA:  e non includere altro testo nella tua risposta.} \\
\\ 
\texttt{torta, oliva, aquila, fantasma, veleno, bufalo, canada, tazza, colpo, pupilla, rivoluzione, germania, morte, ornitorinco, blocco, spazio, lettiera, acero, rete, carta, berlino, radice, talpa, strega, stadio} \\
            }
        }
    }
    & \\ \\

    \theutterance \stepcounter{utterance}  
    & & & \multicolumn{4}{p{0.6\linewidth}}{
        \cellcolor[rgb]{0.9,0.9,0.9}{
            \makecell[{{p{\linewidth}}}]{
                \texttt{\tiny{[GM$\langle$P2]}}
                \texttt{INDOVINA: canada, germania} \\
            }
        }
    }
     \\ \\

    \theutterance \stepcounter{utterance}  
    & & & \multicolumn{2}{p{0.3\linewidth}}{
        \cellcolor[rgb]{0.9,0.9,0.9}{
            \makecell[{{p{\linewidth}}}]{
                \texttt{\tiny{[GM$|$GM]}}
                ['canada', 'germania']
            }
        }
    }
    & & \\ \\

    \theutterance \stepcounter{utterance}  
    & & & \multicolumn{2}{p{0.3\linewidth}}{
        \cellcolor[rgb]{0.9,0.9,0.9}{
            \makecell[{{p{\linewidth}}}]{
                \texttt{\tiny{[GM$|$GM]}}
                {'word': 'canada', 'assignment': 'team'}
            }
        }
    }
    & & \\ \\

    \theutterance \stepcounter{utterance}  
    & & & \multicolumn{2}{p{0.3\linewidth}}{
        \cellcolor[rgb]{0.9,0.9,0.9}{
            \makecell[{{p{\linewidth}}}]{
                \texttt{\tiny{[GM$|$GM]}}
                {'word': 'canada', 'assignment': 'team'}
            }
        }
    }
    & & \\ \\

    \theutterance \stepcounter{utterance}  
    & & & \multicolumn{2}{p{0.3\linewidth}}{
        \cellcolor[rgb]{0.9,0.9,0.9}{
            \makecell[{{p{\linewidth}}}]{
                \texttt{\tiny{[GM$|$GM]}}
                {'word': 'germania', 'assignment': 'team'}
            }
        }
    }
    & & \\ \\

    \theutterance \stepcounter{utterance}  
    & & & \multicolumn{2}{p{0.3\linewidth}}{
        \cellcolor[rgb]{0.9,0.9,0.9}{
            \makecell[{{p{\linewidth}}}]{
                \texttt{\tiny{[GM$|$GM]}}
                {'word': 'germania', 'assignment': 'team'}
            }
        }
    }
    & & \\ \\

    \theutterance \stepcounter{utterance}  
    & & & \multicolumn{2}{p{0.3\linewidth}}{
        \cellcolor[rgb]{0.9,0.9,0.9}{
            \makecell[{{p{\linewidth}}}]{
                \texttt{\tiny{[GM$|$GM]}}
                {'word': 'pupilla', 'assignment': 'opponent'}
            }
        }
    }
    & & \\ \\

    \theutterance \stepcounter{utterance}  
    & & & \multicolumn{2}{p{0.3\linewidth}}{
        \cellcolor[rgb]{0.9,0.9,0.9}{
            \makecell[{{p{\linewidth}}}]{
                \texttt{\tiny{[GM$|$GM]}}
                {'hidden': {'team': ['ornitorinco', 'lettiera', 'veleno', 'spazio', 'radice', 'rete', 'tazza'], 'innocent': ['blocco', 'morte', 'aquila'], 'opponent': ['oliva', 'fantasma', 'strega', 'stadio', 'torta', 'acero', 'carta'], 'assassin': ['bufalo', 'colpo', 'rivoluzione', 'berlino', 'talpa']}, 'revealed': {'team': {'team': ['canada', 'germania'], 'innocent': [], 'opponent': [], 'assassin': []}, 'opponent': {'team': [], 'innocent': [], 'opponent': ['pupilla'], 'assassin': []}}}
            }
        }
    }
    & & \\ \\

    \theutterance \stepcounter{utterance}  
    & & \multicolumn{4}{p{0.6\linewidth}}{
        \cellcolor[rgb]{0.9,0.9,0.9}{
            \makecell[{{p{\linewidth}}}]{
                \texttt{\tiny{[P1$\langle$GM]}}
                \texttt{Le parole canada, germania sono state indovinate correttamente. Il turno del tuo compagno di squadra è terminato lì.} \\
\texttt{Ora fornisci un altro indizio relativo ad alcune delle parole rimanenti della tua squadra e un elenco delle parole target correlate. Ricorda di iniziare il tuo indizio con 'INDIZIO: ', inserire una nuova riga e iniziare il tuo elenco di parole target separate da virgole con 'OBIETTIVI: '. Nota: alcune parole sono state rimosse dagli elenchi rispetto alle richieste precedenti.} \\
\\ 
\texttt{Le parole rimanenti della tua squadra sono: ornitorinco, lettiera, veleno, spazio, radice, rete, tazza.} \\
\texttt{Le parole rimanenti del tuo avversario sono: oliva, fantasma, strega, stadio, pupilla, torta, acero, carta.} \\
\texttt{Le parole distrattive rimanenti sono: blocco, morte, aquila.} \\
\texttt{Le parole mortali rimanenti sono: bufalo, colpo, rivoluzione, berlino, talpa.} \\
            }
        }
    }
    & \\ \\

    \theutterance \stepcounter{utterance}  
    & \multicolumn{4}{p{0.6\linewidth}}{
        \cellcolor[rgb]{0.9,0.9,0.9}{
            \makecell[{{p{\linewidth}}}]{
                \texttt{\tiny{[P1$\rangle$GM]}}
                \texttt{INDIZIO: Tossico} \\
\texttt{OBIETTIVI: ornitorinco, veleno} \\
            }
        }
    }
    & & \\ \\

    \theutterance \stepcounter{utterance}  
    & & & \multicolumn{2}{p{0.3\linewidth}}{
        \cellcolor[rgb]{0.9,0.9,0.9}{
            \makecell[{{p{\linewidth}}}]{
                \texttt{\tiny{[GM$|$GM]}}
                \texttt{tossico} \\
            }
        }
    }
    & & \\ \\

    \theutterance \stepcounter{utterance}  
    & & & \multicolumn{2}{p{0.3\linewidth}}{
        \cellcolor[rgb]{0.9,0.9,0.9}{
            \makecell[{{p{\linewidth}}}]{
                \texttt{\tiny{[GM$|$GM]}}
                ['ornitorinco', 'veleno']
            }
        }
    }
    & & \\ \\

    \theutterance \stepcounter{utterance}  
    & & & \multicolumn{2}{p{0.3\linewidth}}{
        \cellcolor[rgb]{0.9,0.9,0.9}{
            \makecell[{{p{\linewidth}}}]{
                \texttt{\tiny{[GM$|$GM]}}
                {'word': 'ornitorinco', 'assignment': 'team'}
            }
        }
    }
    & & \\ \\

    \theutterance \stepcounter{utterance}  
    & & & \multicolumn{2}{p{0.3\linewidth}}{
        \cellcolor[rgb]{0.9,0.9,0.9}{
            \makecell[{{p{\linewidth}}}]{
                \texttt{\tiny{[GM$|$GM]}}
                {'word': 'veleno', 'assignment': 'team'}
            }
        }
    }
    & & \\ \\

    \theutterance \stepcounter{utterance}  
    & & \multicolumn{4}{p{0.6\linewidth}}{
        \cellcolor[rgb]{0.9,0.9,0.9}{
            \makecell[{{p{\linewidth}}}]{
                \texttt{\tiny{[GM$\rangle$P2]}}
                \texttt{Le parole canada, germania sono state indovinate correttamente. Il tuo turno è terminato lì.} \\
\texttt{Ora fornisci un altro elenco separato da virgole di almeno 1 e fino a 2 parole dal seguente elenco di parole che meglio si riferiscono o sono più strettamente associate alla parola 'tossico'. Ricorda di iniziare la tua risposta con INDOVINA: . Nota: alcune parole sono state rimosse dall'elenco rispetto alle richieste precedenti.} \\
\\ 
\texttt{torta, oliva, aquila, fantasma, veleno, bufalo, tazza, colpo, rivoluzione, morte, ornitorinco, blocco, spazio, lettiera, acero, rete, carta, berlino, radice, talpa, strega, stadio} \\
            }
        }
    }
    & \\ \\

    \theutterance \stepcounter{utterance}  
    & & & \multicolumn{4}{p{0.6\linewidth}}{
        \cellcolor[rgb]{0.9,0.9,0.9}{
            \makecell[{{p{\linewidth}}}]{
                \texttt{\tiny{[GM$\langle$P2]}}
                \texttt{INDOVINA: veleno, ornitorinco} \\
            }
        }
    }
     \\ \\

    \theutterance \stepcounter{utterance}  
    & & & \multicolumn{2}{p{0.3\linewidth}}{
        \cellcolor[rgb]{0.9,0.9,0.9}{
            \makecell[{{p{\linewidth}}}]{
                \texttt{\tiny{[GM$|$GM]}}
                ['veleno', 'ornitorinco']
            }
        }
    }
    & & \\ \\

    \theutterance \stepcounter{utterance}  
    & & & \multicolumn{2}{p{0.3\linewidth}}{
        \cellcolor[rgb]{0.9,0.9,0.9}{
            \makecell[{{p{\linewidth}}}]{
                \texttt{\tiny{[GM$|$GM]}}
                {'word': 'veleno', 'assignment': 'team'}
            }
        }
    }
    & & \\ \\

    \theutterance \stepcounter{utterance}  
    & & & \multicolumn{2}{p{0.3\linewidth}}{
        \cellcolor[rgb]{0.9,0.9,0.9}{
            \makecell[{{p{\linewidth}}}]{
                \texttt{\tiny{[GM$|$GM]}}
                {'word': 'veleno', 'assignment': 'team'}
            }
        }
    }
    & & \\ \\

    \theutterance \stepcounter{utterance}  
    & & & \multicolumn{2}{p{0.3\linewidth}}{
        \cellcolor[rgb]{0.9,0.9,0.9}{
            \makecell[{{p{\linewidth}}}]{
                \texttt{\tiny{[GM$|$GM]}}
                {'word': 'ornitorinco', 'assignment': 'team'}
            }
        }
    }
    & & \\ \\

    \theutterance \stepcounter{utterance}  
    & & & \multicolumn{2}{p{0.3\linewidth}}{
        \cellcolor[rgb]{0.9,0.9,0.9}{
            \makecell[{{p{\linewidth}}}]{
                \texttt{\tiny{[GM$|$GM]}}
                {'word': 'ornitorinco', 'assignment': 'team'}
            }
        }
    }
    & & \\ \\

    \theutterance \stepcounter{utterance}  
    & & & \multicolumn{2}{p{0.3\linewidth}}{
        \cellcolor[rgb]{0.9,0.9,0.9}{
            \makecell[{{p{\linewidth}}}]{
                \texttt{\tiny{[GM$|$GM]}}
                {'word': 'carta', 'assignment': 'opponent'}
            }
        }
    }
    & & \\ \\

    \theutterance \stepcounter{utterance}  
    & & & \multicolumn{2}{p{0.3\linewidth}}{
        \cellcolor[rgb]{0.9,0.9,0.9}{
            \makecell[{{p{\linewidth}}}]{
                \texttt{\tiny{[GM$|$GM]}}
                {'hidden': {'team': ['lettiera', 'spazio', 'radice', 'rete', 'tazza'], 'innocent': ['blocco', 'morte', 'aquila'], 'opponent': ['oliva', 'fantasma', 'strega', 'stadio', 'torta', 'acero'], 'assassin': ['bufalo', 'colpo', 'rivoluzione', 'berlino', 'talpa']}, 'revealed': {'team': {'team': ['canada', 'germania', 'veleno', 'ornitorinco'], 'innocent': [], 'opponent': [], 'assassin': []}, 'opponent': {'team': [], 'innocent': [], 'opponent': ['pupilla', 'carta'], 'assassin': []}}}
            }
        }
    }
    & & \\ \\

    \theutterance \stepcounter{utterance}  
    & & \multicolumn{4}{p{0.6\linewidth}}{
        \cellcolor[rgb]{0.9,0.9,0.9}{
            \makecell[{{p{\linewidth}}}]{
                \texttt{\tiny{[P1$\langle$GM]}}
                \texttt{Le parole veleno, ornitorinco sono state indovinate correttamente. Il turno del tuo compagno di squadra è terminato lì.} \\
\texttt{Ora fornisci un altro indizio relativo ad alcune delle parole rimanenti della tua squadra e un elenco delle parole target correlate. Ricorda di iniziare il tuo indizio con 'INDIZIO: ', inserire una nuova riga e iniziare il tuo elenco di parole target separate da virgole con 'OBIETTIVI: '. Nota: alcune parole sono state rimosse dagli elenchi rispetto alle richieste precedenti.} \\
\\ 
\texttt{Le parole rimanenti della tua squadra sono: lettiera, spazio, radice, rete, tazza.} \\
\texttt{Le parole rimanenti del tuo avversario sono: oliva, fantasma, strega, stadio, torta, acero, carta.} \\
\texttt{Le parole distrattive rimanenti sono: blocco, morte, aquila.} \\
\texttt{Le parole mortali rimanenti sono: bufalo, colpo, rivoluzione, berlino, talpa.} \\
            }
        }
    }
    & \\ \\

    \theutterance \stepcounter{utterance}  
    & \multicolumn{4}{p{0.6\linewidth}}{
        \cellcolor[rgb]{0.9,0.9,0.9}{
            \makecell[{{p{\linewidth}}}]{
                \texttt{\tiny{[P1$\rangle$GM]}}
                \texttt{INDIZIO: Web} \\
\texttt{OBIETTIVI: rete, spazio} \\
            }
        }
    }
    & & \\ \\

    \theutterance \stepcounter{utterance}  
    & & & \multicolumn{2}{p{0.3\linewidth}}{
        \cellcolor[rgb]{0.9,0.9,0.9}{
            \makecell[{{p{\linewidth}}}]{
                \texttt{\tiny{[GM$|$GM]}}
                \texttt{web} \\
            }
        }
    }
    & & \\ \\

    \theutterance \stepcounter{utterance}  
    & & & \multicolumn{2}{p{0.3\linewidth}}{
        \cellcolor[rgb]{0.9,0.9,0.9}{
            \makecell[{{p{\linewidth}}}]{
                \texttt{\tiny{[GM$|$GM]}}
                ['rete', 'spazio']
            }
        }
    }
    & & \\ \\

    \theutterance \stepcounter{utterance}  
    & & & \multicolumn{2}{p{0.3\linewidth}}{
        \cellcolor[rgb]{0.9,0.9,0.9}{
            \makecell[{{p{\linewidth}}}]{
                \texttt{\tiny{[GM$|$GM]}}
                {'word': 'rete', 'assignment': 'team'}
            }
        }
    }
    & & \\ \\

    \theutterance \stepcounter{utterance}  
    & & & \multicolumn{2}{p{0.3\linewidth}}{
        \cellcolor[rgb]{0.9,0.9,0.9}{
            \makecell[{{p{\linewidth}}}]{
                \texttt{\tiny{[GM$|$GM]}}
                {'word': 'spazio', 'assignment': 'team'}
            }
        }
    }
    & & \\ \\

    \theutterance \stepcounter{utterance}  
    & & \multicolumn{4}{p{0.6\linewidth}}{
        \cellcolor[rgb]{0.9,0.9,0.9}{
            \makecell[{{p{\linewidth}}}]{
                \texttt{\tiny{[GM$\rangle$P2]}}
                \texttt{Le parole veleno, ornitorinco sono state indovinate correttamente. Il tuo turno è terminato lì.} \\
\texttt{Ora fornisci un altro elenco separato da virgole di almeno 1 e fino a 2 parole dal seguente elenco di parole che meglio si riferiscono o sono più strettamente associate alla parola 'web'. Ricorda di iniziare la tua risposta con INDOVINA: . Nota: alcune parole sono state rimosse dall'elenco rispetto alle richieste precedenti.} \\
\\ 
\texttt{torta, oliva, aquila, fantasma, bufalo, tazza, colpo, rivoluzione, morte, blocco, spazio, lettiera, acero, rete, berlino, radice, talpa, strega, stadio} \\
            }
        }
    }
    & \\ \\

    \theutterance \stepcounter{utterance}  
    & & & \multicolumn{4}{p{0.6\linewidth}}{
        \cellcolor[rgb]{0.9,0.9,0.9}{
            \makecell[{{p{\linewidth}}}]{
                \texttt{\tiny{[GM$\langle$P2]}}
                \texttt{INDOVINA: rete, strega} \\
            }
        }
    }
     \\ \\

    \theutterance \stepcounter{utterance}  
    & & & \multicolumn{2}{p{0.3\linewidth}}{
        \cellcolor[rgb]{0.9,0.9,0.9}{
            \makecell[{{p{\linewidth}}}]{
                \texttt{\tiny{[GM$|$GM]}}
                ['rete', 'strega']
            }
        }
    }
    & & \\ \\

    \theutterance \stepcounter{utterance}  
    & & & \multicolumn{2}{p{0.3\linewidth}}{
        \cellcolor[rgb]{0.9,0.9,0.9}{
            \makecell[{{p{\linewidth}}}]{
                \texttt{\tiny{[GM$|$GM]}}
                {'word': 'rete', 'assignment': 'team'}
            }
        }
    }
    & & \\ \\

    \theutterance \stepcounter{utterance}  
    & & & \multicolumn{2}{p{0.3\linewidth}}{
        \cellcolor[rgb]{0.9,0.9,0.9}{
            \makecell[{{p{\linewidth}}}]{
                \texttt{\tiny{[GM$|$GM]}}
                {'word': 'rete', 'assignment': 'team'}
            }
        }
    }
    & & \\ \\

    \theutterance \stepcounter{utterance}  
    & & & \multicolumn{2}{p{0.3\linewidth}}{
        \cellcolor[rgb]{0.9,0.9,0.9}{
            \makecell[{{p{\linewidth}}}]{
                \texttt{\tiny{[GM$|$GM]}}
                {'word': 'strega', 'assignment': 'opponent'}
            }
        }
    }
    & & \\ \\

    \theutterance \stepcounter{utterance}  
    & & & \multicolumn{2}{p{0.3\linewidth}}{
        \cellcolor[rgb]{0.9,0.9,0.9}{
            \makecell[{{p{\linewidth}}}]{
                \texttt{\tiny{[GM$|$GM]}}
                \texttt{strega} \\
            }
        }
    }
    & & \\ \\

    \theutterance \stepcounter{utterance}  
    & & & \multicolumn{2}{p{0.3\linewidth}}{
        \cellcolor[rgb]{0.9,0.9,0.9}{
            \makecell[{{p{\linewidth}}}]{
                \texttt{\tiny{[GM$|$GM]}}
                {'word': 'oliva', 'assignment': 'opponent'}
            }
        }
    }
    & & \\ \\

    \theutterance \stepcounter{utterance}  
    & & & \multicolumn{2}{p{0.3\linewidth}}{
        \cellcolor[rgb]{0.9,0.9,0.9}{
            \makecell[{{p{\linewidth}}}]{
                \texttt{\tiny{[GM$|$GM]}}
                {'hidden': {'team': ['lettiera', 'spazio', 'radice', 'tazza'], 'innocent': ['blocco', 'morte', 'aquila'], 'opponent': ['fantasma', 'stadio', 'torta', 'acero'], 'assassin': ['bufalo', 'colpo', 'rivoluzione', 'berlino', 'talpa']}, 'revealed': {'team': {'team': ['canada', 'germania', 'veleno', 'ornitorinco', 'rete'], 'innocent': [], 'opponent': ['strega'], 'assassin': []}, 'opponent': {'team': [], 'innocent': [], 'opponent': ['pupilla', 'carta', 'oliva'], 'assassin': []}}}
            }
        }
    }
    & & \\ \\

    \theutterance \stepcounter{utterance}  
    & & \multicolumn{4}{p{0.6\linewidth}}{
        \cellcolor[rgb]{0.9,0.9,0.9}{
            \makecell[{{p{\linewidth}}}]{
                \texttt{\tiny{[P1$\langle$GM]}}
                \texttt{La parola rete è stata indovinata correttamente. La parola strega è stata indovinata, ma è una parola avversario. Il turno del tuo compagno di squadra è terminato lì.} \\
\texttt{Ora fornisci un altro indizio relativo ad alcune delle parole rimanenti della tua squadra e un elenco delle parole target correlate. Ricorda di iniziare il tuo indizio con 'INDIZIO: ', inserire una nuova riga e iniziare il tuo elenco di parole target separate da virgole con 'OBIETTIVI: '. Nota: alcune parole sono state rimosse dagli elenchi rispetto alle richieste precedenti.} \\
\\ 
\texttt{Le parole rimanenti della tua squadra sono: lettiera, spazio, radice, tazza.} \\
\texttt{Le parole rimanenti del tuo avversario sono: oliva, fantasma, stadio, torta, acero.} \\
\texttt{Le parole distrattive rimanenti sono: blocco, morte, aquila.} \\
\texttt{Le parole mortali rimanenti sono: bufalo, colpo, rivoluzione, berlino, talpa.} \\
            }
        }
    }
    & \\ \\

    \theutterance \stepcounter{utterance}  
    & \multicolumn{4}{p{0.6\linewidth}}{
        \cellcolor[rgb]{0.9,0.9,0.9}{
            \makecell[{{p{\linewidth}}}]{
                \texttt{\tiny{[P1$\rangle$GM]}}
                \texttt{INDIZIO: Bagno} \\
\texttt{OBIETTIVI: lettiera, tazza} \\
            }
        }
    }
    & & \\ \\

    \theutterance \stepcounter{utterance}  
    & & & \multicolumn{2}{p{0.3\linewidth}}{
        \cellcolor[rgb]{0.9,0.9,0.9}{
            \makecell[{{p{\linewidth}}}]{
                \texttt{\tiny{[GM$|$GM]}}
                \texttt{bagno} \\
            }
        }
    }
    & & \\ \\

    \theutterance \stepcounter{utterance}  
    & & & \multicolumn{2}{p{0.3\linewidth}}{
        \cellcolor[rgb]{0.9,0.9,0.9}{
            \makecell[{{p{\linewidth}}}]{
                \texttt{\tiny{[GM$|$GM]}}
                ['lettiera', 'tazza']
            }
        }
    }
    & & \\ \\

    \theutterance \stepcounter{utterance}  
    & & & \multicolumn{2}{p{0.3\linewidth}}{
        \cellcolor[rgb]{0.9,0.9,0.9}{
            \makecell[{{p{\linewidth}}}]{
                \texttt{\tiny{[GM$|$GM]}}
                {'word': 'lettiera', 'assignment': 'team'}
            }
        }
    }
    & & \\ \\

    \theutterance \stepcounter{utterance}  
    & & & \multicolumn{2}{p{0.3\linewidth}}{
        \cellcolor[rgb]{0.9,0.9,0.9}{
            \makecell[{{p{\linewidth}}}]{
                \texttt{\tiny{[GM$|$GM]}}
                {'word': 'tazza', 'assignment': 'team'}
            }
        }
    }
    & & \\ \\

    \theutterance \stepcounter{utterance}  
    & & \multicolumn{4}{p{0.6\linewidth}}{
        \cellcolor[rgb]{0.9,0.9,0.9}{
            \makecell[{{p{\linewidth}}}]{
                \texttt{\tiny{[GM$\rangle$P2]}}
                \texttt{La parola rete è stata indovinata correttamente. La parola strega è stata indovinata, ma è una parola avversario. Il tuo turno è terminato lì.} \\
\texttt{Ora fornisci un altro elenco separato da virgole di almeno 1 e fino a 2 parole dal seguente elenco di parole che meglio si riferiscono o sono più strettamente associate alla parola 'bagno'. Ricorda di iniziare la tua risposta con INDOVINA: . Nota: alcune parole sono state rimosse dall'elenco rispetto alle richieste precedenti.} \\
\\ 
\texttt{torta, aquila, fantasma, bufalo, tazza, colpo, rivoluzione, morte, blocco, spazio, lettiera, acero, berlino, radice, talpa, stadio} \\
            }
        }
    }
    & \\ \\

    \theutterance \stepcounter{utterance}  
    & & & \multicolumn{4}{p{0.6\linewidth}}{
        \cellcolor[rgb]{0.9,0.9,0.9}{
            \makecell[{{p{\linewidth}}}]{
                \texttt{\tiny{[GM$\langle$P2]}}
                \texttt{INDOVINA: tazza, lettiera} \\
            }
        }
    }
     \\ \\

    \theutterance \stepcounter{utterance}  
    & & & \multicolumn{2}{p{0.3\linewidth}}{
        \cellcolor[rgb]{0.9,0.9,0.9}{
            \makecell[{{p{\linewidth}}}]{
                \texttt{\tiny{[GM$|$GM]}}
                ['tazza', 'lettiera']
            }
        }
    }
    & & \\ \\

    \theutterance \stepcounter{utterance}  
    & & & \multicolumn{2}{p{0.3\linewidth}}{
        \cellcolor[rgb]{0.9,0.9,0.9}{
            \makecell[{{p{\linewidth}}}]{
                \texttt{\tiny{[GM$|$GM]}}
                {'word': 'tazza', 'assignment': 'team'}
            }
        }
    }
    & & \\ \\

    \theutterance \stepcounter{utterance}  
    & & & \multicolumn{2}{p{0.3\linewidth}}{
        \cellcolor[rgb]{0.9,0.9,0.9}{
            \makecell[{{p{\linewidth}}}]{
                \texttt{\tiny{[GM$|$GM]}}
                {'word': 'tazza', 'assignment': 'team'}
            }
        }
    }
    & & \\ \\

    \theutterance \stepcounter{utterance}  
    & & & \multicolumn{2}{p{0.3\linewidth}}{
        \cellcolor[rgb]{0.9,0.9,0.9}{
            \makecell[{{p{\linewidth}}}]{
                \texttt{\tiny{[GM$|$GM]}}
                {'word': 'lettiera', 'assignment': 'team'}
            }
        }
    }
    & & \\ \\

    \theutterance \stepcounter{utterance}  
    & & & \multicolumn{2}{p{0.3\linewidth}}{
        \cellcolor[rgb]{0.9,0.9,0.9}{
            \makecell[{{p{\linewidth}}}]{
                \texttt{\tiny{[GM$|$GM]}}
                {'word': 'lettiera', 'assignment': 'team'}
            }
        }
    }
    & & \\ \\

    \theutterance \stepcounter{utterance}  
    & & & \multicolumn{2}{p{0.3\linewidth}}{
        \cellcolor[rgb]{0.9,0.9,0.9}{
            \makecell[{{p{\linewidth}}}]{
                \texttt{\tiny{[GM$|$GM]}}
                {'word': 'acero', 'assignment': 'opponent'}
            }
        }
    }
    & & \\ \\

    \theutterance \stepcounter{utterance}  
    & & & \multicolumn{2}{p{0.3\linewidth}}{
        \cellcolor[rgb]{0.9,0.9,0.9}{
            \makecell[{{p{\linewidth}}}]{
                \texttt{\tiny{[GM$|$GM]}}
                {'hidden': {'team': ['spazio', 'radice'], 'innocent': ['blocco', 'morte', 'aquila'], 'opponent': ['fantasma', 'stadio', 'torta'], 'assassin': ['bufalo', 'colpo', 'rivoluzione', 'berlino', 'talpa']}, 'revealed': {'team': {'team': ['canada', 'germania', 'veleno', 'ornitorinco', 'rete', 'tazza', 'lettiera'], 'innocent': [], 'opponent': ['strega'], 'assassin': []}, 'opponent': {'team': [], 'innocent': [], 'opponent': ['pupilla', 'carta', 'oliva', 'acero'], 'assassin': []}}}
            }
        }
    }
    & & \\ \\

    \theutterance \stepcounter{utterance}  
    & & \multicolumn{4}{p{0.6\linewidth}}{
        \cellcolor[rgb]{0.9,0.9,0.9}{
            \makecell[{{p{\linewidth}}}]{
                \texttt{\tiny{[P1$\langle$GM]}}
                \texttt{Le parole tazza, lettiera sono state indovinate correttamente. Il turno del tuo compagno di squadra è terminato lì.} \\
\texttt{Ora fornisci un altro indizio relativo ad alcune delle parole rimanenti della tua squadra e un elenco delle parole target correlate. Ricorda di iniziare il tuo indizio con 'INDIZIO: ', inserire una nuova riga e iniziare il tuo elenco di parole target separate da virgole con 'OBIETTIVI: '. Nota: alcune parole sono state rimosse dagli elenchi rispetto alle richieste precedenti.} \\
\\ 
\texttt{Le parole rimanenti della tua squadra sono: spazio, radice.} \\
\texttt{Le parole rimanenti del tuo avversario sono: fantasma, stadio, torta, acero.} \\
\texttt{Le parole distrattive rimanenti sono: blocco, morte, aquila.} \\
\texttt{Le parole mortali rimanenti sono: bufalo, colpo, rivoluzione, berlino, talpa.} \\
            }
        }
    }
    & \\ \\

    \theutterance \stepcounter{utterance}  
    & \multicolumn{4}{p{0.6\linewidth}}{
        \cellcolor[rgb]{0.9,0.9,0.9}{
            \makecell[{{p{\linewidth}}}]{
                \texttt{\tiny{[P1$\rangle$GM]}}
                \texttt{INDIZIO: Cosmo} \\
\texttt{OBIETTIVI: spazio, radice} \\
            }
        }
    }
    & & \\ \\

    \theutterance \stepcounter{utterance}  
    & & & \multicolumn{2}{p{0.3\linewidth}}{
        \cellcolor[rgb]{0.9,0.9,0.9}{
            \makecell[{{p{\linewidth}}}]{
                \texttt{\tiny{[GM$|$GM]}}
                \texttt{cosmo} \\
            }
        }
    }
    & & \\ \\

    \theutterance \stepcounter{utterance}  
    & & & \multicolumn{2}{p{0.3\linewidth}}{
        \cellcolor[rgb]{0.9,0.9,0.9}{
            \makecell[{{p{\linewidth}}}]{
                \texttt{\tiny{[GM$|$GM]}}
                ['spazio', 'radice']
            }
        }
    }
    & & \\ \\

    \theutterance \stepcounter{utterance}  
    & & & \multicolumn{2}{p{0.3\linewidth}}{
        \cellcolor[rgb]{0.9,0.9,0.9}{
            \makecell[{{p{\linewidth}}}]{
                \texttt{\tiny{[GM$|$GM]}}
                {'word': 'spazio', 'assignment': 'team'}
            }
        }
    }
    & & \\ \\

    \theutterance \stepcounter{utterance}  
    & & & \multicolumn{2}{p{0.3\linewidth}}{
        \cellcolor[rgb]{0.9,0.9,0.9}{
            \makecell[{{p{\linewidth}}}]{
                \texttt{\tiny{[GM$|$GM]}}
                {'word': 'radice', 'assignment': 'team'}
            }
        }
    }
    & & \\ \\

    \theutterance \stepcounter{utterance}  
    & & \multicolumn{4}{p{0.6\linewidth}}{
        \cellcolor[rgb]{0.9,0.9,0.9}{
            \makecell[{{p{\linewidth}}}]{
                \texttt{\tiny{[GM$\rangle$P2]}}
                \texttt{Le parole tazza, lettiera sono state indovinate correttamente. Il tuo turno è terminato lì.} \\
\texttt{Ora fornisci un altro elenco separato da virgole di almeno 1 e fino a 2 parole dal seguente elenco di parole che meglio si riferiscono o sono più strettamente associate alla parola 'cosmo'. Ricorda di iniziare la tua risposta con INDOVINA: . Nota: alcune parole sono state rimosse dall'elenco rispetto alle richieste precedenti.} \\
\\ 
\texttt{torta, aquila, fantasma, bufalo, colpo, rivoluzione, morte, blocco, spazio, berlino, radice, talpa, stadio} \\
            }
        }
    }
    & \\ \\

    \theutterance \stepcounter{utterance}  
    & & & \multicolumn{4}{p{0.6\linewidth}}{
        \cellcolor[rgb]{0.9,0.9,0.9}{
            \makecell[{{p{\linewidth}}}]{
                \texttt{\tiny{[GM$\langle$P2]}}
                \texttt{INDOVINA: spazio} \\
            }
        }
    }
     \\ \\

    \theutterance \stepcounter{utterance}  
    & & & \multicolumn{2}{p{0.3\linewidth}}{
        \cellcolor[rgb]{0.9,0.9,0.9}{
            \makecell[{{p{\linewidth}}}]{
                \texttt{\tiny{[GM$|$GM]}}
                ['spazio']
            }
        }
    }
    & & \\ \\

    \theutterance \stepcounter{utterance}  
    & & & \multicolumn{2}{p{0.3\linewidth}}{
        \cellcolor[rgb]{0.9,0.9,0.9}{
            \makecell[{{p{\linewidth}}}]{
                \texttt{\tiny{[GM$|$GM]}}
                {'word': 'spazio', 'assignment': 'team'}
            }
        }
    }
    & & \\ \\

    \theutterance \stepcounter{utterance}  
    & & & \multicolumn{2}{p{0.3\linewidth}}{
        \cellcolor[rgb]{0.9,0.9,0.9}{
            \makecell[{{p{\linewidth}}}]{
                \texttt{\tiny{[GM$|$GM]}}
                {'word': 'spazio', 'assignment': 'team'}
            }
        }
    }
    & & \\ \\

    \theutterance \stepcounter{utterance}  
    & & & \multicolumn{2}{p{0.3\linewidth}}{
        \cellcolor[rgb]{0.9,0.9,0.9}{
            \makecell[{{p{\linewidth}}}]{
                \texttt{\tiny{[GM$|$GM]}}
                {'word': 'torta', 'assignment': 'opponent'}
            }
        }
    }
    & & \\ \\

    \theutterance \stepcounter{utterance}  
    & & & \multicolumn{2}{p{0.3\linewidth}}{
        \cellcolor[rgb]{0.9,0.9,0.9}{
            \makecell[{{p{\linewidth}}}]{
                \texttt{\tiny{[GM$|$GM]}}
                {'hidden': {'team': ['radice'], 'innocent': ['blocco', 'morte', 'aquila'], 'opponent': ['fantasma', 'stadio'], 'assassin': ['bufalo', 'colpo', 'rivoluzione', 'berlino', 'talpa']}, 'revealed': {'team': {'team': ['canada', 'germania', 'veleno', 'ornitorinco', 'rete', 'tazza', 'lettiera', 'spazio'], 'innocent': [], 'opponent': ['strega'], 'assassin': []}, 'opponent': {'team': [], 'innocent': [], 'opponent': ['pupilla', 'carta', 'oliva', 'acero', 'torta'], 'assassin': []}}}
            }
        }
    }
    & & \\ \\

    \theutterance \stepcounter{utterance}  
    & & \multicolumn{4}{p{0.6\linewidth}}{
        \cellcolor[rgb]{0.9,0.9,0.9}{
            \makecell[{{p{\linewidth}}}]{
                \texttt{\tiny{[P1$\langle$GM]}}
                \texttt{La parola spazio è stata indovinata correttamente. Il turno del tuo compagno di squadra è terminato lì.} \\
\texttt{Ora fornisci un altro indizio relativo ad alcune delle parole rimanenti della tua squadra e un elenco delle parole target correlate. Ricorda di iniziare il tuo indizio con 'INDIZIO: ', inserire una nuova riga e iniziare il tuo elenco di parole target separate da virgole con 'OBIETTIVI: '. Nota: alcune parole sono state rimosse dagli elenchi rispetto alle richieste precedenti.} \\
\\ 
\texttt{Le parole rimanenti della tua squadra sono: radice.} \\
\texttt{Le parole rimanenti del tuo avversario sono: fantasma, stadio, torta.} \\
\texttt{Le parole distrattive rimanenti sono: blocco, morte, aquila.} \\
\texttt{Le parole mortali rimanenti sono: bufalo, colpo, rivoluzione, berlino, talpa.} \\
            }
        }
    }
    & \\ \\

    \theutterance \stepcounter{utterance}  
    & \multicolumn{4}{p{0.6\linewidth}}{
        \cellcolor[rgb]{0.9,0.9,0.9}{
            \makecell[{{p{\linewidth}}}]{
                \texttt{\tiny{[P1$\rangle$GM]}}
                \texttt{INDIZIO: Albero} \\
\texttt{OBIETTIVI: radice} \\
            }
        }
    }
    & & \\ \\

    \theutterance \stepcounter{utterance}  
    & & & \multicolumn{2}{p{0.3\linewidth}}{
        \cellcolor[rgb]{0.9,0.9,0.9}{
            \makecell[{{p{\linewidth}}}]{
                \texttt{\tiny{[GM$|$GM]}}
                \texttt{albero} \\
            }
        }
    }
    & & \\ \\

    \theutterance \stepcounter{utterance}  
    & & & \multicolumn{2}{p{0.3\linewidth}}{
        \cellcolor[rgb]{0.9,0.9,0.9}{
            \makecell[{{p{\linewidth}}}]{
                \texttt{\tiny{[GM$|$GM]}}
                ['radice']
            }
        }
    }
    & & \\ \\

    \theutterance \stepcounter{utterance}  
    & & & \multicolumn{2}{p{0.3\linewidth}}{
        \cellcolor[rgb]{0.9,0.9,0.9}{
            \makecell[{{p{\linewidth}}}]{
                \texttt{\tiny{[GM$|$GM]}}
                {'word': 'radice', 'assignment': 'team'}
            }
        }
    }
    & & \\ \\

    \theutterance \stepcounter{utterance}  
    & & \multicolumn{4}{p{0.6\linewidth}}{
        \cellcolor[rgb]{0.9,0.9,0.9}{
            \makecell[{{p{\linewidth}}}]{
                \texttt{\tiny{[GM$\rangle$P2]}}
                \texttt{La parola spazio è stata indovinata correttamente. Il tuo turno è terminato lì.} \\
\texttt{Ora fornisci un altro elenco separato da virgole di almeno 1 e fino a 1 parole dal seguente elenco di parole che meglio si riferiscono o sono più strettamente associate alla parola 'albero'. Ricorda di iniziare la tua risposta con INDOVINA: . Nota: alcune parole sono state rimosse dall'elenco rispetto alle richieste precedenti.} \\
\\ 
\texttt{aquila, fantasma, bufalo, colpo, rivoluzione, morte, blocco, berlino, radice, talpa, stadio} \\
            }
        }
    }
    & \\ \\

    \theutterance \stepcounter{utterance}  
    & & & \multicolumn{4}{p{0.6\linewidth}}{
        \cellcolor[rgb]{0.9,0.9,0.9}{
            \makecell[{{p{\linewidth}}}]{
                \texttt{\tiny{[GM$\langle$P2]}}
                \texttt{INDOVINA: radice} \\
            }
        }
    }
     \\ \\

    \theutterance \stepcounter{utterance}  
    & & & \multicolumn{2}{p{0.3\linewidth}}{
        \cellcolor[rgb]{0.9,0.9,0.9}{
            \makecell[{{p{\linewidth}}}]{
                \texttt{\tiny{[GM$|$GM]}}
                ['radice']
            }
        }
    }
    & & \\ \\

    \theutterance \stepcounter{utterance}  
    & & & \multicolumn{2}{p{0.3\linewidth}}{
        \cellcolor[rgb]{0.9,0.9,0.9}{
            \makecell[{{p{\linewidth}}}]{
                \texttt{\tiny{[GM$|$GM]}}
                {'word': 'radice', 'assignment': 'team'}
            }
        }
    }
    & & \\ \\

    \theutterance \stepcounter{utterance}  
    & & & \multicolumn{2}{p{0.3\linewidth}}{
        \cellcolor[rgb]{0.9,0.9,0.9}{
            \makecell[{{p{\linewidth}}}]{
                \texttt{\tiny{[GM$|$GM]}}
                {'word': 'radice', 'assignment': 'team'}
            }
        }
    }
    & & \\ \\

    \theutterance \stepcounter{utterance}  
    & & & \multicolumn{2}{p{0.3\linewidth}}{
        \cellcolor[rgb]{0.9,0.9,0.9}{
            \makecell[{{p{\linewidth}}}]{
                \texttt{\tiny{[GM$|$GM]}}
                \texttt{team has won} \\
            }
        }
    }
    & & \\ \\

\end{supertabular}
}

\end{document}
